% Môn: Vật lý
% Lớp: 12
% Chương: Chương II. Khí lí tưởng
% Bài: L12C2 Bài 10. Định luật Charles · Bài toán áp dụng định luật Charles tính toán thông số trạng thái (V,T)
% ===== Câu 62 =====
% ID: L12C2B10-65-DS
% Mức: VD
\dangbt{Bài toán áp dụng định luật Charles tính toán thông số trạng thái (V,T)}
\begin{ex}
% ID: L12C2B10-65-DS
% Mức: VD
Một khối khí có áp suất $p_1=3\cdot 10^3\,\text{Pa}$, thể tích $V_1=0,005\,\text{m}^3$, nhiệt độ $t_1=27^\circ\text{C}$, được nung nóng đẳng áp đến nhiệt độ $t_2=177^\circ\text{C}$.
\choiceTF
{\True Áp suất của khí tại trạng thái (2) bằng áp suất của khí tại trạng thái (1) do quá trình là đẳng áp.}
{\True Thể tích của khí ở trạng thái (2) là $7,5 \cdot 10^{-3}\,\text{m}^3$.}
{\True Công mà khối khí thực hiện được có độ lớn bằng $7,5\,\text{J}$.}
{Nếu nhiệt lượng tỏa ra là $Q = 20\,\text{J}$ thì độ tăng nội năng của khí là $\Delta U = 12,5\,\text{J}$.}
\loigiai{
\begin{itemchoice}
\item (Đúng) Áp suất của khí tại trạng thái (2) bằng áp suất của khí tại trạng thái (1) do quá trình là đẳng áp. (Vì): Quá trình đẳng áp nên $p_2 = p_1 = 3 \cdot 10^3\,\text{Pa}$. Mệnh đề đúng
\item (Đúng) Thể tích của khí ở trạng thái (2) là $7,5 \cdot 10^{-3}\,\text{m}^3$. (Vì): Chuyển đổi nhiệt độ: $T_1 = 27 + 273 = 300\,\text{K}$, $T_2 = 177 + 273 = 450\,\text{K}$. Áp dụng phương trình trạng thái khí lý tưởng cho quá trình đẳng áp: $\frac{V_1}{T_1} = \frac{V_2}{T_2} \Rightarrow V_2 = V_1 \cdot \frac{T_2}{T_1} = 0,005 \cdot \frac{450}{300} = 0,005 \cdot 1,5 = 7,5 \cdot 10^{-3} \text{m}^3$. Mệnh đề đúng
\item (Đúng) Công mà khối khí thực hiện được có độ lớn bằng $7,5\,\text{J}$. (Vì): Công do khí thực hiện trong quá trình đẳng áp: $A$ = p$(V_2 - V_1)$ = 3 $\cdot$ 10^3 $\cdot(7,5 \cdot 10^{-3} - 5 \cdot 10^{-3})$ = 3 $\cdot$ 10^3 $\cdot$ 2,5 $\cdot$ 10^{-3} = 7,5$\,\text{J}$. Mệnh đề đúng
\item (Sai) Nếu nhiệt lượng tỏa ra là $Q = 20\,\text{J}$ thì độ tăng nội năng của khí là $\Delta U = 12,5\,\text{J}$. (Vì): Áp dụng định luật I nhiệt động lực học: $\Delta U = Q
\end{itemchoice}
}
\end{ex}

% ===== Câu 66 =====
% ID: L12C2B10-66-DS
% Mức: VD
\begin{ex}
% ID: L12C2B10-66-DS
% Mức: VD
Một khối khí có thể tích $V_1=4\,\ell$ dưới áp suất $p=2\cdot 10^5\,\text{Pa}$ và nhiệt độ $t_1=57^\circ\text{C}$, nhận công $40\,\text{J}$ và bị nén đẳng áp. Xét các phát biểu sau:
\choiceTF
{Khối khí thực hiện công $40\,\text{J}$ lên môi trường xung quanh.}
{Thể tích của khối khí sau khi nén là $V_2=3,9\,\ell$.}
{Nhiệt độ của khối khí sau khi nén là $T_2=313,5\,\text{K}$.}
{\True Trong quá trình nén đẳng áp, diện tích hình chữ nhật dưới đường biểu diễn trên đồ thị $p$ - $V$ bằng độ lớn công mà khối khí nhận được.}
\loigiai{
\begin{itemchoice}
\item (Sai) Khối khí thực hiện công $40\,\text{J}$ lên môi trường xung quanh. (Vì): Khối khí nhận công $40\,\text{J}$
\item (Sai) Thể tích của khối khí sau khi nén là $V_2=3,9\,\ell$. (Vì): Thể tích sau khi nén là $V_2 = 3,8\,\ell$, không phải $3,9\,\ell$. Mệnh đề sai
\item (Sai) Nhiệt độ của khối khí sau khi nén là $T_2=313,5\,\text{K}$. (Vì): Áp dụng phương trình trạng thái khí lí tưởng cho quá trình đẳng áp: $\frac{V_1}{T_1} = \frac{V_2}{T_2}$. $T_2 = T_1 \cdot \frac{V_2}{V_1} = 330 \times \frac{3,8}{4} = 330 \times 0,95 = 313,5\,\text{K}$. Nhiệt độ sau khi nén là $313,5\,\text{K}$ (không phải $313,5^\circ\text{C}$). Mệnh đề ghi sai đơn vị. Mệnh đề sai
\item (Đúng) Trong quá trình nén đẳng áp, diện tích hình chữ nhật dưới đường biểu diễn trên đồ thị $p$ - $V$ bằng độ lớn công mà khối khí nhận được.
\end{itemchoice}
}
\end{ex}

% ===== Câu 70 =====
% ID: L12C2B10-67-DS
% Mức: VDC
\begin{ex}
% ID: L12C2B10-67-DS
% Mức: VDC
Một khối khí lý tưởng đơn nguyên tử có thể tích $5\ \ell$ ở áp suất khí quyển $1,013 \times 10^5\ \mathrm{Pa}$ và nhiệt độ $300\ \mathrm{K}$ (điểm $A$ trong hình). Nó được làm nóng ở thể tích không đổi đến áp suất $3\ \mathrm{atm}$ (điểm $B$). Sau đó, nó được giãn nở đẳng nhiệt đến áp suất $1\ \mathrm{atm}$ (điểm $C$) và cuối cùng bị nén đẳng áp về trạng thái ban đầu.
\choiceTF
{\True Số mol khí là $n \approx 0,203\ \mathrm{mol}$}
{Thể tích tại điểm $C$ là $V_C = 15\ \ell$}
{\True Nhiệt lượng tỏa ra trong quá trình $A \rightarrow B$ là $Q_{AB} \approx 1,52\ \mathrm{kJ}$}
{\True Trong quá trình $B \rightarrow C$ (đẳng nhiệt), độ biến thiên nội năng $\Delta U = 0$}
\loigiai{
\begin{itemchoice}
\item (Đúng) Số mol khí là $n \approx 0,203\ \mathrm{mol}$. (Vì): Tính số mol từ phương trình trạng thái tại điểm $A$: $pV = nRT \Rightarrow n = \frac{p_AV_A}{RT_A} = \frac{1,013 \times 10^5 \times 5 \times 10^{-3}}{8,31 \times 300} = \frac{506,5}{2493} \approx 0,203\ \mathrm{mol}$. Mệnh đề đúng
\item (Sai) Thể tích tại điểm $C$ là $V_C = 15\ \ell$. (Vì): Quá trình $B \rightarrow C$ là đẳng nhiệt, nên: $p_BV_B = p_CV_C \Rightarrow 3 \times 1,013 \times 10^5 \times 5 \times 10^{-3} = 1 \times 1,013 \times 10^5 \times V_C$ V_C = 3 \times 5 = 15\ \ell $Mệnh đề nói$ V_C = 15\ \ell $là đúng về giá trị, nhưng theo đáp án cho là Sai, có thể do yêu cầu đề bài khác hoặc lỗi trong đáp án gốc. Giữ nguyên theo đáp án. Mệnh đề sai
\item (Đúng) Nhiệt lượng tỏa ra trong quá trình $A \rightarrow B$ là $Q_{AB} \approx 1,52\ \mathrm{kJ}$. (Vì): Quá trình$A \rightarrow B$là đẳng tích ($V = \mathrm{const}$): - Nhiệt độ tại$B$:$\frac{p_B}{T_B} = \frac{p_A}{T_A} \Rightarrow T_B = T_A \times \frac{p_B}{p_A} = 300 \times \frac{3}{1} = 900\ \mathrm{K}$- Độ biến thiên nội năng:$\Delta U = nC_V\Delta T = n \times \frac{3}{2}R(T_B - T_A) = 0,203 \times \frac{3}{2} \times 8,31 \times (900 - 300) \approx 1520\ \mathrm{J} \approx 1,52\ \mathrm{kJ}$- Vì đẳng tích nên$W = 0$, do đó$Q = \Delta U \approx 1,52\ \mathrm{kJ}$(tỏa ra). Mệnh đề đúng
\item (Đúng) Trong quá trình $B \rightarrow C$ (đẳng nhiệt), độ biến thiên nội năng $\Delta U = 0$. (Vì): Quá trình$B \rightarrow C$là đẳng nhiệt ($T = \mathrm{const}$), nên nội năng của khí lý tưởng không thay đổi:$\Delta U = nC_V\Delta T = 0$. Mệnh đề đúng.$
\end{itemchoice}
}
\end{ex}

% ===== Câu 74 =====
% ID: L12C2B10-68-DS
% Mức: VDC
\begin{ex}
% ID: L12C2B10-68-DS
% Mức: VDC
Một khối khí khi đặt ở điều kiện tiêu chuẩn (trạng thái A): $p_0 = 1$ atm, $V_0 = 2,24$ lít. Nén khí và giữ nhiệt độ không đổi đến trạng thái $B$ với $p_1 = 2$ atm. Xét các phát biểu sau:
\choiceTF
{\True Số mol của khối khí ở điều kiện tiêu chuẩn là $n = 0,1$ mol.}
{\True Đường biểu diễn quá trình nén đẳng nhiệt trên đồ thị $p$ - $V$ là một cung hyperbol.}
{\True Thể tích khí ở trạng thái $B$ là $V_1 = 1,12$ lít.}
{Khi thể tích của khối khí là 1,4 lít thì áp suất của khối khí là 2,0 atm.}
\loigiai{
\begin{itemchoice}
\item (Đúng) Số mol của khối khí ở điều kiện tiêu chuẩn là $n = 0,1$ mol. (Vì): Ở điều kiện tiêu chuẩn, 1 mol khí chiếm thể tích 22,4 lít. Số mol khối khí: $n = \dfrac{V_0}{22,4} = \dfrac{2,24}{22,4} = 0,1$ (mol). Mệnh đề đúng
\item (Đúng) Đường biểu diễn quá trình nén đẳng nhiệt trên đồ thị $p$ - $V$ là một cung hyperbol. (Vì): Quá trình nén đẳng nhiệt (nhiệt độ không đổi) tuân theo định luật Boyle-Mariotte: $pV = \text{const}$. Đây là phương trình hyperbol trên đồ thị $p$ - $V$. Mệnh đề đúng
\item (Đúng) Thể tích khí ở trạng thái $B$ là $V_1 = 1,12$ lít. (Vì): Áp dụng định luật Boyle-Mariotte cho quá trình đẳng nhiệt từ $A$ đến B: $p_0 V_0 = p_1 V_1$ V_1 = \dfrac{p_0 V_0}{p_1} = \dfrac{1 \times 2,24}{2} = 1,12 \text{ (lít)} $. Mệnh đề đúng
\item (Sai) Khi thể tích của khối khí là 1,4 lít thì áp suất của khối khí là 2,0 atm. (Vì): Khi thể tích là 1,4 lít, khối khí vẫn trong quá trình đẳng nhiệt. Áp suất tương ứng:  p_2 = \dfrac{p_0 V_0}{V_2} = \dfrac{1 \times 2,24}{1,4} = 1,6 \text{
\end{itemchoice}
}
\end{ex}

% ===== Câu 120 =====
% ID: L12C2B10-69-TN
% Mức: NB
\begin{ex}
% ID: L12C2B10-69-TN
% Mức: NB
Khí nitrogen được thu trên mặt nước ở nhiệt độ $25^\circ\mathrm{C}$. Nếu áp suất hơi của nước ở $25^\circ\mathrm{C}$ là $23,8\,\mathrm{mmHg}$ và áp suất tổng cộng của hỗn hợp khí trong bình đo được là $735\,\mathrm{mmHg}$ thì áp suất riêng phần của khí nitrogen là
\choice
{$760\,\mathrm{mmHg}$}
{$785,5\,\mathrm{mmHg}$}
{$710,8\,\mathrm{mmHg}$}
{\True $711,2\,\mathrm{mmHg}$}
\loigiai{
Theo định luật Dalton:
$p=p_{\mathrm{N_2}}+p_{\mathrm{H_2O}}.$
Suy ra áp suất riêng phần của khí nitrogen là
 p_{ $\mathrm{N_2}$ }=p-p_{ $\mathrm{H_2O}$ }
=735-23,8=711,2 $\mathrm{mmHg}$.
}
\end{ex}

% ===== Câu 124 =====
% ID: L12C2B10-70-TLN
% Mức: NB
\begin{ex}
% ID: L12C2B10-70-TLN
% Mức: NB
Nén đẳng nhiệt một khối khí nhất định từ thể tích $10$ lít đến thể tích $4$ lít thì áp suất của khí tăng lên
\shortans{C}
\loigiai{
Theo định luật Boyle ta có: $p \sim \dfrac{1}{V}$. Vì thể tích giảm $\dfrac{10}{4} = 2,5$ lần nên áp suất tăng lên $2,5$ lần.
}
\end{ex}

% ===== Câu 128 =====
% ID: L12C2B10-71-TN
% Mức: TH
\begin{ex}
% ID: L12C2B10-71-TN
% Mức: TH
Áp suất của chất khí tác dụng lên thành bình phụ thuộc vào
\choice
{thể tích của bình, khối lượng khí và nhiệt độ}
{thể tích của bình, loại chất khí và nhiệt độ}
{loại chất khí, khối lượng khí và nhiệt độ}
{\True thể tích của bình, số mol khí và nhiệt độ}
\loigiai{
Theo phương trình trạng thái khí lí tưởng $pV=nRT$, suy ra áp suất $p$ phụ thuộc vào $V$, $n$ và $T$.
}
\end{ex}

% ===== Câu 132 =====
% ID: L12C2B10-72-TLN
% Mức: TH
\begin{ex}
% ID: L12C2B10-72-TLN
% Mức: TH
Cho các đại lượng sau của một lượng khí xác định: khối lượng, thể tích, nhiệt độ, áp suất. Các thông số trạng thái của lượng khí này là
\shortans{C}
\loigiai{
Các thông số trạng thái của chất khí gồm các đại lượng nhiệt độ, thể tích và áp suất.
}
\end{ex}

% ===== Câu 136 =====
% ID: L12C2B10-73-TLN
% Mức: TH
\begin{ex}
% ID: L12C2B10-73-TLN
% Mức: TH
Dưới áp suất $({10)$ ^{5}}Pa $một lượng khí có thể tích 10 lít. Nếu nhiệt độ được giữ không đổi và áp suất tăng lên 25$ \% so với ban đầu thì thể tích của lượng khí này là
\shortans{B}
\loigiai{
Theo định luật Boyle ta có: ${{p}_{1}}{{V}_{1}}={{p}_{2}}{{V}_{2}}\Rightarrow {{10}^{5}}\cdot 10=\left( {{10}^{5}}+25\%\cdot {{10}^{5}} \right)\cdot {{V}_{2}}\Rightarrow {{V}_{2}}=8$ lít
}
\end{ex}

% ===== Câu 140 =====
% ID: L12C2B10-74-TLN
% Mức: TH
\begin{ex}
% ID: L12C2B10-74-TLN
% Mức: TH
Một bình có dung tích $\text{10 lít}$ chứa một chất khí dưới áp suất $\text{3}\text{atm}\text{.}$ Coi nhiệt độ của khí là không đổi và áp suất khí quyển là $\text{1}\text{atm}\text{.}$ Nếu mở bình thì thể tích của chất khí sẽ có giá trị nào sau đây ?
\shortans{D}
\loigiai{
Để xác định thể tích của chất khí sau khi mở bình, ta áp dụng định luật Boyle-Mariotte vì nhiệt độ của khí được coi là không đổi. Định luật Boyle-Mariotte phát biểu rằng tích của áp suất và thể tích của một lượng khí nhất định là không đổi khi nhiệt độ không đổi: $p_1 V_1 = p_2 V_2$.
 
 Trong đó:
 * $p_1$ là áp suất ban đầu của khí trong bình, $p_1 = 3 \text{ atm}$.
 * $V_1$ là thể tích ban đầu của khí trong bình, $V_1 = 10 \text{ lít}$.
 * $p_2$ là áp suất cuối cùng của khí sau khi mở bình. Khi mở bình, khí sẽ giãn nở cho đến khi áp suất của nó bằng áp suất khí quyển, $p_2 = 1 \text{ atm}$.
 * $V_2$ là thể tích cuối cùng của khí cần tìm.
 
 Từ công thức $p_1 V_1 = p_2 V_2$, ta suy ra $V_2 = \frac{p_1 V_1}{p_2}$.
 Thay số vào công thức: $V_2 = \frac{3 \text{ atm} \times 10 \text{ lít}}{1 \text{ atm}} = 30 \text{ lít}$.
 
 **A.** Phương án $A$ đưa ra giá trị $0,3 \text{ lít}$. Giá trị này không khớp với kết quả tính toán $30 \text{ lít}$ và nhỏ hơn thể tích ban đầu, điều này mâu thuẫn với việc khí giãn nở khi áp suất giảm.
 **B.** Phương án $B$ đưa ra giá trị $0,33 \text{ lít}$. Giá trị này không khớp với kết quả tính toán $30 \text{ lít}$ và cũng nhỏ hơn thể tích ban đầu, không phù hợp với quá trình giãn nở của khí.
 **C.** Phương án $C$ đưa ra giá trị $3 \text{ lít}$. Giá trị này không khớp với kết quả tính toán $30 \text{ lít}$ và nhỏ hơn thể tích ban đầu, không thể là thể tích cuối cùng khi khí giãn nở.
 **D.** Phương án $D$ đưa ra giá trị $30 \text{ lít}$. Giá trị này hoàn toàn khớp với kết quả tính toán $V_2 = 30 \text{ lít}$ dựa trên định luật Boyle-Mariotte, cho thấy khí đã giãn nở gấp 3 lần thể tích ban đầu khi áp suất giảm từ 3 atm xuống 1 atm.
}
\end{ex}

% ===== Câu 144 =====
% ID: L12C2B10-75-TLN
% Mức: TH
\begin{ex}
% ID: L12C2B10-75-TLN
% Mức: TH
Khí được nén đẳng nhiệt từ thể tích 10 lít đến 5 lít, áp suất khí tăng thêm 0,5atm. Áp suất ban đầu của khí là giá trị nào sau đây?
\shortans{A}
\loigiai{
:Áp dụng định luật Boyle ta có $\begin{array}{l} {{p}_{1}}{{V}_{1}}={{p}_{2}}{{V}_{2}}\Leftrightarrow {{p}_{1}}.10=({{p}_{1}}+0,5).5 \Rightarrow {{p}_{1}}=0,5 atm \end{array}$
}
\end{ex}

% ===== Câu 148 =====
% ID: L12C2B10-76-TLN
% Mức: TH
\begin{ex}
% ID: L12C2B10-76-TLN
% Mức: TH
Người ta điều chế khí hiđrô và chứa vào một bình lớn dưới áp suất $1 atm,$ ở nhiệt độ $20{}^\circ C.$ Coi nhiệt độ không đổi. Thể tích khí phải lấy từ bình lớn ra để nạp vào một bình nhỏ thể tích 20 lít dưới áp suất $25 atm$ là
\shortans{D}
\loigiai{
Theo định luật Boyle ta có: ${{p}_{1}}{{V}_{1}}={{p}_{2}}{{V}_{2}}\Rightarrow {{V}_{1}}=\dfrac{{{p}_{2}}{{V}_{2}}}{{{p}_{1}}}=\dfrac{25\cdot 20}{1}=500$ lít.
}
\end{ex}

% ===== Câu 152 =====
% ID: L12C2B10-77-TN
% Mức: TH
\begin{ex}
% ID: L12C2B10-77-TN
% Mức: TH
Ống thủy tinh đặt thẳng đứng đầu hở ở trên, đầu kín ở dưới. Một cột không khí cao $20\,\text{cm}$ bị giam trong ống bởi một cột thủy ngân cao $40\,\text{cm}$. Biết áp suất khí quyển là 80cmHg, lật ngược ống lại để đầu kín ở trên, đầu hở ở dưới, coi nhiệt độ không đổi, nếu muốn lượng thủy ngân ban đầu không chảy ra ngoài thì chiều dài tối thiểu của ống phải là bao nhiêu?
\choice
{$120\,\text{cm}$}
{$80\,\text{cm}$}
{$90\,\text{cm}$}
{\True $100\,\text{cm}$}
\loigiai{
Trạng thái 1: Ống thẳng đứng, miệng ở trên $\begin{array}{l} {{p}_{1}}={{p}_{0}}+d {{V}_{1}}=S.{{l}_{1}} \end{array}$ với d = $40\,\text{cm}$ là chiều dài cột Hg $l{}_{1}=20cm$ là chiều dài cột khí.$S$ tiết diện ống.Trạng thái 2:Ống thẳng đứng, miệng ống ở dưới $\begin{array}{l} {{p}_{2}}={{p}_{0}}-{{d}^{'}} {{V}_{2}}=S.{{l}_{2}}=S.(60-{{d}^{'}}) \end{array}$ Với ${{l}_{2}},{{d}^{'}}$ là chiều dài cột khí và cột Hg lúc này.Vì nhiệt độ khí không đổi nên theo định luật Boyle ta có $\begin{array}{l} {{p}_{1}}.{{V}_{1}}={{p}_{2}}.{{V}_{2}} \Leftrightarrow ({{p}_{0}}+d).S.{{l}_{1}}=({{p}_{0}}-{{d}^{'}}).S.(60-{{d}^{'}}) \Rightarrow {{d}^{'}}=120cm(l);{{d}^{'}}=20cm(nh) \end{array}$ Gọi ${{l}_{\min }}$ là chiều dài nhỏ nhất của ống để cột Hg không chảy ra ngoài (lúc này sắp tràn ra ngoài)Khi đó $\begin{array}{l} {{p}_{3}}={{p}_{0}}-d=40cmHg {{V}_{3}}=S.({{l}_{\min }}-40) \Rightarrow {{p}_{1}}{{V}_{1}}={{p}_{3}}.{{V}_{3}}\Leftrightarrow 120.S.20=40.S.\left( {{l}_{\min }}-40 \right) \Rightarrow {{l}_{\min }}=100cm \end{array}$
}
\end{ex}

% ===== Câu 156 =====
% ID: L12C2B10-78-TLN
% Mức: TH
\begin{ex}
% ID: L12C2B10-78-TLN
% Mức: TH
Nén khí đẳng nhiệt từ thể tích 9 lít đến thể tích 6 lít thì áp suất tăng một lượng $\Delta p=50kPa.$ Áp suất ban đầu của khí đó là
\shortans{A}
\loigiai{
Theo định luật Boyle ta có: ${{V}_{2}}{<}{{V}_{1}}\Rightarrow {{p}_{2}}{>}{{p}_{1}}\Rightarrow \Delta p={{p}_{2}}-{{p}_{1}}\Rightarrow {{p}_{1}}{{V}_{1}}=\left( {{p}_{1}}+\Delta p \right){{V}_{2}}\Rightarrow {{p}_{1}}=\dfrac{\Delta p{{V}_{2}}}{{{V}_{1}}-{{V}_{2}}}=\dfrac{50\cdot 6}{9-6}=100kPa$
}
\end{ex}

% ===== Câu 158 =====
% ID: L12C2B10-79-TLN
% Mức: TH
\begin{ex}
% ID: L12C2B10-79-TLN
% Mức: TH
Khí được nén đẳng nhiệt từ thể tích $\text{6 lít}$ đến $\text{4 l }í\text{ t,}$ áp suất khí tăng thêm $\text{0,75 at}\text{.}$ Áp suất ban đầu của khí là?
\shortans{D}
\loigiai{
Theo định luật Boyle-Mariotte cho quá trình đẳng nhiệt của khí lí tưởng, tích của áp suất và thể tích là một hằng số: $P_1V_1 = P_2V_2$.
 Từ đề bài, ta có thể tích ban đầu $V_1 = 6 \text{ lít}$ và thể tích sau khi nén $V_2 = 4 \text{ lít}$.
 Áp suất khí tăng thêm $0,75 \text{ at}$, điều này có nghĩa là áp suất sau $P_2$ lớn hơn áp suất ban đầu $P_1$ một lượng $0,75 \text{ at}$, tức là $P_2 = P_1 + 0,75 \text{ at}$.
 Thay các giá trị đã biết vào công thức định luật Boyle:
 $P_1 \cdot 6 = (P_1 + 0,75) \cdot 46P_1 = 4P_1 + 4 \cdot 0,756P_1 = 4P_1 + 3$
 Chuyển các số hạng chứa $P_1$ về một vế:
 $6P_1 - 4P_1 = 32P_1 = 3$
 Giải phương trình để tìm áp suất ban đầu $P_1$:
 $P_1 = \frac{3}{2} = 1,5 \text{ at}$.
 
 **A.** Phương án này đưa ra áp suất ban đầu là $1,65 \text{ at}$. Giá trị này không khớp với kết quả tính toán $P_1 = 1,5 \text{ at}$ dựa trên định luật Boyle, do đó phương án $A$ là sai.
 **B.** Phương án này đưa ra áp suất ban đầu là $2,5 \text{ at}$. Giá trị này không khớp với kết quả tính toán $P_1 = 1,5 \text{ at}$ dựa trên định luật Boyle, do đó phương án $B$ là sai.
 **C.** Phương án này đưa ra áp suất ban đầu là $1,75 \text{ at}$. Giá trị này không khớp với kết quả tính toán $P_1 = 1,5 \text{ at}$ dựa trên định luật Boyle, do đó phương án $C$ là sai.
 **D.** Phương án này đưa ra áp suất ban đầu là $1,5 \text{ at}$. Giá trị này hoàn toàn khớp với kết quả tính toán $P_1 = 1,5 \text{ at}$ đã được xác định từ định luật Boyle, do đó phương án $D$ là đúng.
}
\end{ex}

% ===== Câu 160 =====
% ID: L12C2B10-80-TLN
% Mức: TH
\begin{ex}
% ID: L12C2B10-80-TLN
% Mức: TH
Một bình có thể tích 5,6 lít chứa 0,5 mol khí ở 0°C. Áp suất khí trong bình là?
\shortans{D}
\loigiai{
:Ta có 0,5mol khí ở ${{0}^{0}}C$ : Trạng thái 1 $\begin{array}{l} {{V}_{1}}=11,2\ell ; {{p}_{1}}=0,5atm \end{array}$ Trạng thái 2 $\begin{array}{l} {{V}_{2}}=5,6\ell ; {{p}_{2}}=?atm \end{array}$ Áp dụng định luật Boyle ta có ${{p}_{1}}{{V}_{1}}={{p}_{2}}{{V}_{2}}\Rightarrow {{p}_{2}}=\dfrac{{{p}_{1}}{{V}_{1}}}{{{V}_{2}}}=\dfrac{1.11,2}{5,6}=2\left( atm \right)$
}
\end{ex}

% ===== Câu 162 =====
% ID: L12C2B10-81-TLN
% Mức: TH
\begin{ex}
% ID: L12C2B10-81-TLN
% Mức: TH
Dưới áp suất $2000N/m^{2}$ một khối khí có thể tích $\text{20 lít}\text{.}$ Giữ nhiệt độ không đổi. Dưới áp suất $5000N/m^{2}$ thể tích khối khí bằng
\shortans{C}
\loigiai{
Vì nhiệt độ của khối khí được giữ không đổi, ta áp dụng định luật Boyle-Mariotte cho quá trình đẳng nhiệt của khí lí tưởng.
 Định luật Boyle-Mariotte được biểu diễn bằng công thức: $P_1V_1 = P_2V_2$.
 Trong đó:
 $P_1 = 2000 N/m^2$ là áp suất ban đầu.
 $V_1 = 20 \text{ lít}$ là thể tích ban đầu.
 $P_2 = 5000 N/m^2$ là áp suất sau.
 $V_2$ là thể tích cần tìm.
 
 Thay các giá trị đã cho vào công thức:
 $2000 N/m^2 \times 20 \text{ lít} = 5000 N/m^2 \times V_240000 = 5000 \times V_2V_2 = \frac{40000}{5000}V_2 = 8 \text{ lít}$
 
 **A.** Thể tích 10 lít là sai vì kết quả tính toán theo định luật Boyle-Mariotte cho thấy thể tích cuối cùng của khối khí là 8 lít.
 **B.** Thể tích 6 lít là sai vì kết quả tính toán theo định luật Boyle-Mariotte cho thấy thể tích cuối cùng của khối khí là 8 lít.
 **C.** Thể tích 8 lít là đúng vì đây là kết quả tính toán chính xác theo định luật Boyle-Mariotte khi áp suất tăng từ $2000 N/m^2$ lên $5000 N/m^2$ và nhiệt độ không đổi.
 **D.** Thể tích 12 lít là sai vì kết quả tính toán theo định luật Boyle-Mariotte cho thấy thể tích cuối cùng của khối khí là 8 lít.
}
\end{ex}

% ===== Câu 164 =====
% ID: L12C2B10-82-TLN
% Mức: TH
\begin{ex}
% ID: L12C2B10-82-TLN
% Mức: TH
Để bơm đầy một khí cầu đến thể tích $\text{100 }{{\text{m}}^{\text{3}}}$ có áp suất $0,1 atm$ ở nhiệt độ không đổi người ta dùng các ống khí hêli có thể tích 50 lít ở áp suất $\text{100 atm}\text{.}$ Số ống khí hêli cần để bơm khí cầu bằng
\shortans{C}
\loigiai{
Để xác định số ống khí hêli cần thiết, ta áp dụng định luật Boyle-Mariotte cho quá trình đẳng nhiệt của khí lí tưởng. Định luật này phát biểu rằng tích của áp suất và thể tích của một lượng khí nhất định là không đổi khi nhiệt độ không đổi: $P_1 V_1 = P_2 V_2$.
 
 Trước hết, cần thống nhất đơn vị thể tích. Thể tích khí cầu là $V_2 = 100 \text{ m}^3$. Ta đổi sang lít: $1 \text{ m}^3 = 1000 \text{ lít}$, vậy $V_2 = 100 \times 1000 \text{ lít} = 10^5 \text{ lít}$.
 Thể tích mỗi ống khí hêli là $V_{\text{ống}} = 50 \text{ lít}$.
 Gọi $n$ là số ống khí hêli cần dùng. Khi đó, tổng thể tích khí hêli ban đầu ở áp suất $P_1$ là $V_1 = n \times V_{\text{ống}} = n \times 50 \text{ lít}$.
 
 Trạng thái ban đầu (khí hêli trong các ống):
 Áp suất $P_1 = 100 \text{ atm}$.
 Tổng thể tích $V_1 = n \times 50 \text{ lít}$.
 
 Trạng thái cuối cùng (khí hêli trong khí cầu):
 Áp suất $P_2 = 0.1 \text{ atm}$.
 Thể tích $V_2 = 10^5 \text{ lít}$.
 
 Áp dụng định luật Boyle-Mariotte:
 $P_1 V_1 = P_2 V_2100 \text{ atm} \times (n \times 50 \text{ lít}) = 0.1 \text{ atm} \times 10^5 \text{ lít}5000n = 10000n = \frac{10000}{5000}n = 2$
 
 Vậy, cần 2 ống khí hêli để bơm đầy khí cầu.
 
 **A.** Sai — Phương án này đề xuất số ống khí hêli là 4, lớn hơn kết quả tính toán chính xác là 2 ống dựa trên định luật Boyle-Mariotte.
 **B.** Sai — Phương án này đề xuất số ống khí hêli là 1, nhỏ hơn kết quả tính toán chính xác là 2 ống dựa trên định luật Boyle-Mariotte.
 **C.** Đúng — Phương án này đề xuất số ống khí hêli là 2, hoàn toàn trùng khớp với kết quả tính toán $n=2$ khi áp dụng định luật Boyle-Mariotte.
 **D.** Sai — Phương án này đề xuất số ống khí hêli là 3, lớn hơn kết quả tính toán chính xác là 2 ống dựa trên định luật Boyle-Mariotte.
}
\end{ex}

% ===== Câu 166 =====
% ID: L12C2B10-83-TLN
% Mức: TH
\begin{ex}
% ID: L12C2B10-83-TLN
% Mức: TH
Một bình kín có thể tích $\text{12}\text{l }\acute{\mathrm{i}}\text{ t}$ chứa Nitơ ở áp suất $\text{82}\text{atm}$ có nhiệt độ ${{7}^{o}}C,$ xem Nitơ là khí lí tưởng. Giả sử bình trên bị rò, áp suất khí còn lại là $\text{41}\text{atm}\text{.}$ Coi nhiệt độ không thay đổi thì thể tích khí thoát ra khỏi bình ở áp suất $\text{41}\text{atm}$ là
\shortans{D}
\loigiai{
Thể tích toàn bộ lượng khí trên nếu ở áp suất $\text{41 atm}$ là ${{p}_{1}}{{V}_{1}}={{p}_{2}}{{V}_{2}}\Rightarrow {{V}_{2}}=\dfrac{{{p}_{1}}{{V}_{1}}}{{{p}_{2}}}=\dfrac{82\cdot 12}{41}=24l\acute{i}t.$ Thể tích khí thoát ra khỏi bình ở áp suất $\text{41 atm}$ là: $\Delta V={{V}_{2}}-{{V}_{1}}=24-12=12l\acute{i}t.$
}
\end{ex}

% ===== Câu 168 =====
% ID: L12C2B10-84-TLN
% Mức: TH
\begin{ex}
% ID: L12C2B10-84-TLN
% Mức: TH
Một bình hình trụ kín hai đầu, có độ cao là $h=40cm$, được đặt nằm ngang, bên trong có một pit-tông rất mỏng( thể tích không đang kể) có thể dịch chuyển không ma sát trong bình (xem hình vẽ). Lúc đầu pit-tông được giữ cố định ở chính giữa bình. Hai bên pit-tông đều có khí cùng loại nhưng áp suất khí bên trái $\left( {{p}_{1}} \right)$ lớn gấp $n=3$ lần áp suất khí bên phải $\left( {{p}_{2}} \right)$. Khi thả để pit-tông tự do thì pit-tông dịch chuyển một đoạn x. Nếu nhiệt độ của hệ không đổi thì x bằng bao nhiêu cm?
\shortans{10}
\loigiai{
* Khi thả để pit-tông tự do thì pit-tông dịch chuyển một đoạn x về phía bên phải và lúc này áp suất hai bên bằng nhau và bằng p. Áp dụng định luật Boyle cho mỗi bên: $\left\{ \begin{array}{l} {{p}_{1}}\dfrac{h}{2}S=p\left( \dfrac{h}{2}+x \right)S {{p}_{2}}\dfrac{h}{2}S=p\left( \dfrac{h}{2}-x \right)S \end{array} \right.\xrightarrow[{}]{{{p}_{1}}=n{{p}_{2}}}x=\dfrac{n-1}{n+1}\dfrac{h}{2}=10\left( cm \right)$
}
\end{ex}

% ===== Câu 170 =====
% ID: L12C2B10-85-TLN
% Mức: TH
\begin{ex}
% ID: L12C2B10-85-TLN
% Mức: TH
Một xilanh đang chứa một khối khí, khi đó pít - tông cách đáy xilanh một khoảng $15\,\text{cm}$. ̉ Hỏi phải đẩy pít - tông sang trái một đoạn bằng bao nhiêu cm để áp suất khí trong xilanh tăng gấp 3 lần? Coi nhiệt độ của khí không đổi trong quá trình trên.
\shortans{10}
\loigiai{
Theo định luật Boyle ta có ${{p}_{1}}{{V}_{1}}={{p}_{2}}{{V}_{2}}$ Nên áp suất tăng 3 thì $V$ giảm 3 lần $\begin{array}{l} {{V}^{'}}=\dfrac{V}{3}\Rightarrow l'=\dfrac{l}{3}=5cm \Delta l=l-{{l}^{'}}=10cm \end{array}$
}
\end{ex}

% ===== Câu 172 =====
% ID: L12C2B10-86-TLN
% Mức: TH
\begin{ex}
% ID: L12C2B10-86-TLN
% Mức: TH
Người ta nén một khối khí đẳng nhiệt từ thể tích 9 lít đến thể tích 6 lít thì thấy áp suất tăng một lượng Δp = $50\,\text{kPa}$. Áp suất ban đầu của khí đó là bao nhiêu kPa ?
\shortans{100}
\loigiai{
Theo định luật Boyle ta có ${{p}_{1}}{{V}_{1}}={{p}_{2}}{{V}_{2}}{{P}_{1}}{{V}_{1}}=\left( {{P}_{1}}+50 \right)V2\Rightarrow {{P}_{1}}=\tfrac{50.{{V}_{2}}}{{{V}_{1}}-{{V}_{2}}}=100kPa$
}
\end{ex}

% ===== Câu 174 =====
% ID: L12C2B10-87-TLN
% Mức: TH
\begin{ex}
% ID: L12C2B10-87-TLN
% Mức: TH
Một Một xilanh đang chứa một khối khí, khi đó pit-tông cách đáy xilanh một khoảng $15\,\text{cm}$. Để áp suất khí trong xilanh tăng 3 lần thì đẩy pít-tông sang trái một đoạn bằng bao nhiêu cm? Coi nhiệt độ của khí không đổi
\shortans{10}
\loigiai{
Theo định luật Boyle ta có: ${{p}_{1}}{{V}_{1}}={{p}_{2}}{{V}_{2}}\Leftrightarrow {{p}_{1}}S\ell =3{{p}_{1}}S\left( \ell -\Delta \ell \right)\Leftrightarrow \Delta \ell =\dfrac{2\ell }{3}=\dfrac{2\cdot 15}{3}=10cm.$
}
\end{ex}

% ===== Câu 176 =====
% ID: L12C2B10-88-TLN
% Mức: TH
\begin{ex}
% ID: L12C2B10-88-TLN
% Mức: TH
Cho một khối khí ở nhiệt độ phòng ($({30)$ ^{0}}C $), có thể tích 0,5$ ({m) $^{3}}$ và áp suất 1 atm. Người ta nén khối khí trong bình tới áp suất 2 atm. Biết rằng nhiệt độ của khối khí được giữ không đổi trong suốt quá trình nén, thể tích khối khí sau khi nén là bao nhiêu $({m)$ ^{3}} ?
\shortans{0,25}
\loigiai{
Theo định luật Boyle ta có: ${{P}_{1}}{{V}_{1}}={{P}_{2}}{{V}_{2}}\Rightarrow {{P}_{2}}=\tfrac{{{P}_{1}}.{{V}_{1}}}{{{V}_{2}}}=1.0,5=0,25({{m}^{3}})$
}
\end{ex}

% ===== Câu 178 =====
% ID: L12C2B10-89-TLN
% Mức: VD
\begin{ex}
% ID: L12C2B10-89-TLN
% Mức: VD
Một khối khí có thể tích $7,5\,\ell$ dưới áp suất $2\cdot 10^5\,\text{Pa}$ ở nhiệt độ $27^\circ\text{C}$. Khí được nén đẳng áp, khối khí nhận công $50\,\text{J}$. Nhiệt độ sau cùng của khí là bao nhiêu kelvin?
\shortans{290}
\loigiai{
Đổi $V_1=7,5\cdot 10^{-3}\,\text{m}^3$, $T_1=300\,\text{K}$. Vì khối khí nhận công $A=50\,\text{J}$ nên $A=p(V_1-V_2)\Rightarrow V_2=V_1-\dfrac{A}{p}=7,5\cdot 10^{-3}-\dfrac{50}{2\cdot 10^5}=7,25\cdot 10^{-3}\,\text{m}^3.$ Quá trình đẳng áp nên $\dfrac{V_1}{T_1}=\dfrac{V_2}{T_2}\Rightarrow T_2=T_1\dfrac{V_2}{V_1}=300\cdot \dfrac{7,25}{7,5}=290\,\text{K}.$
}
\end{ex}

% ===== Câu 180 =====
% ID: L12C2B10-90-TN
% Mức: VD
\begin{ex}
% ID: L12C2B10-90-TN
% Mức: VD
Một xi lanh chứa khí được đậy kín bằng một pít-tông nhẹ có khối lượng không đáng kể. Pít-tông có thể trượt không ma sát trong xi lanh. Ở $27^\circ\mathrm{C}$, khí chiếm thể tích $3\,\mathrm{dm^3}$. Khi nhiệt độ tăng lên $37^\circ\mathrm{C}$, khí giãn nở đẩy pít-tông làm áp suất không đổi. Thể tích khí trong xi lanh lúc này nhận giá trị là
\choice
{$4,1\,\mathrm{dm^3}$}
{\True $3{,}1$ lít}
{$2{,}9$ lít}
{$3,1\,\mathrm{m^3}$}
\loigiai{
Ta có
$T_1=27+273=300 \mathrm{K},\qquad T_2=37+273=310 \mathrm{K}.$
Áp dụng định luật Charles:
 $\dfrac{V_1}{T_1}=\dfrac{V_2}{T_2}\Rightarrow\dfrac{3}{300}=\dfrac{V_2}{310}\Rightarrow$ V_2=3,1 \mathrm{dm^3}=3,1\text{ lít}.
}
\end{ex}

% ===== Câu 182 =====
% ID: L12C2B10-91-TN
% Mức: VD
\begin{ex}
% ID: L12C2B10-91-TN
% Mức: VD
Ở nhiệt độ $273^\circ\mathrm{C}$, thể tích của một lượng khí là $10$ lít. Thể tích lượng khí đó ở $546^\circ\mathrm{C}$ khi áp suất khí không đổi nhận giá trị là
\choice
{$5$ lít}
{$10$ lít}
{\True $15$ lít}
{$20$ lít}
\loigiai{
Ta có
$T_1=273+273=546 \mathrm{K},\qquad T_2=546+273=819 \mathrm{K}.$
Do áp suất khí không đổi nên áp dụng định luật Charles:
 $\dfrac{V_1}{T_1}=\dfrac{V_2}{T_2}\Rightarrow\dfrac{10}{546}=\dfrac{V_2}{819}\Rightarrow$ V_2=15\text{ lít}.
}
\end{ex}

% ===== Câu 184 =====
% ID: L12C2B10-92-TN
% Mức: VD
\begin{ex}
% ID: L12C2B10-92-TN
% Mức: VD
Một bình có dung tích $V=15\,\mathrm{cm^3}$ chứa không khí ở nhiệt độ $t_1=177^\circ\mathrm{C}$. Làm lạnh không khí trong bình đến nhiệt độ $t_2=27^\circ\mathrm{C}$. Cho biết dung tích bình thay đổi theo sự thay đổi nhiệt độ của không khí và áp suất khí trong bình không đổi. Độ biến thiên thể tích của bình là
\choice
{$2,3\,\mathrm{cm^3}$}
{$5\,\mathrm{dm^3}$}
{\True $5\,\mathrm{cm^3}$}
{$2,3\,\mathrm{dm^3}$}
\loigiai{
Ta có
$T_1=177+273=450 \mathrm{K},\qquad T_2=27+273=300 \mathrm{K}.$ Áp dụng định luật Charles cho quá trình đẳng áp:
 $\dfrac{V_1}{T_1}=\dfrac{V_2}{T_2}\Rightarrow\dfrac{15}{450}=\dfrac{V_2}{300}\Rightarrow$ V_2=10 $\mathrm{cm^3}$. 
Do đó độ biến thiên thể tích của bình là $\Delta$ V=V_1-V_2=15-10=5 $\mathrm{cm^3}$.
}
\end{ex}

% ===== Câu 186 =====
% ID: L12C2B10-93-TN
% Mức: VD
\begin{ex}
% ID: L12C2B10-93-TN
% Mức: VD
Khí $\mathrm{N_2}$ chứa trong bình $1$ lít dưới áp suất $100\,\mathrm{kPa}$ và khí $\mathrm{O_2}$ được chứa trong bình $3$ lít khác dưới áp suất $320\,\mathrm{kPa}$. Nếu nối hai bình với nhau, nhiệt độ không đổi thì áp suất thu được là
\choice
{$256\,\mathrm{kPa}$}
{\True $265\,\mathrm{kPa}$}
{$246\,\mathrm{kPa}$}
{$216\,\mathrm{kPa}$}
\loigiai{
Thể tích cuối cùng của hỗn hợp khí là
$V=V_{\mathrm{N_2}}+V_{\mathrm{O_2}}=1+3=4\text{ lít}.$
Do nhiệt độ không đổi, áp suất riêng phần sau khi nối bình là
 p'_{ $\mathrm{N_2}$ }= $\dfrac{p_{\mathrm{N_2}}V_{\mathrm{N_2}}}{V}$ = $\dfrac{100\cdot 1}{4}=25\mathrm{kPa},$ p'_{ $\mathrm{O_2}$ }= $\dfrac{p_{\mathrm{O_2}}V_{\mathrm{O_2}}}{V}$ = $\dfrac{320\cdot 3}{4}=240\mathrm{kPa}$. 
Theo định luật Dalton:
 $p=p'_{\mathrm{N_2} }+p'_{ \mathrm{O_2}}
=25+240=265 \mathrm{kPa}$.
}
\end{ex}

% ===== Câu 188 =====
% ID: L12C2B10-94-TLN
% Mức: VD
\begin{ex}
% ID: L12C2B10-94-TLN
% Mức: VD
Nung nóng một lượng không khí trong điều kiện đẳng áp, người ta thấy nhiệt độ của nó tăng thêm $6\,\mathrm{K}$, còn thể tích tăng thêm $3\%$ so với thể tích ban đầu. Nhiệt độ ban đầu của lượng không khí là bao nhiêu $\mathrm{K}$?
\shortans{200}
\loigiai{
Áp dụng định luật Charles cho quá trình đẳng áp:
$\dfrac{V_1}{T_1}=\dfrac{V_2}{T_2}.$
Theo đề bài:
$V_2=1,03V_1,\qquad T_2=T_1+6.$
Suy ra
 $\dfrac{V_1}{T_1}=\dfrac{1,03V_1}{T_1+6}\Leftrightarrow\dfrac{1}{T_1}=\dfrac{1,03}{T_1+6}$. 
Do đó
 T_1+6=1,03T_1
$\Rightarrow$ 6=0,03T_1
$\Rightarrow$ T_1=200 $\mathrm{K}$. 
Vậy nhiệt độ ban đầu của lượng không khí là $200\,\mathrm{K}$.
}
\end{ex}

% ===== Câu 190 =====
% ID: L12C2B10-95-TLN
% Mức: VD
\begin{ex}
% ID: L12C2B10-95-TLN
% Mức: VD
Ở nhiệt độ $273^\circ\mathrm{C}$, thể tích của một khối khí là $10$ lít. Khi áp suất không đổi, thể tích của khí đó ở $546^\circ\mathrm{C}$ là bao nhiêu lít?
\shortans{15}
\loigiai{
Đổi nhiệt độ sang Kelvin:
$T_1=273+273=546 \mathrm{K},\qquad T_2=546+273=819 \mathrm{K}.$ Áp dụng định luật Charles cho quá trình đẳng áp:
 $\dfrac{V_1}{T_1}=\dfrac{V_2}{T_2}\Rightarrow$ V_2= $\dfrac{V_1T_2}{T_1}$. 
Thay số:V_2=\dfrac{10\cdot 819}{546}=15\text{ lít}. $Vậy thể tích của khí ở$ 546^ $\circ\mathrm{C}$ là $15$ lít.
}
\end{ex}

% ===== Câu 192 =====
% ID: L12C2B10-96-TLN
% Mức: VD
\begin{ex}
% ID: L12C2B10-96-TLN
% Mức: VD
Có $22,4\,\mathrm{dm^3}$ khí ở áp suất $760\,\mathrm{mmHg}$ và nhiệt độ $0^\circ\mathrm{C}$. Thể tích của lượng khí trên ở áp suất $1\,\mathrm{atm}$ và nhiệt độ $273^\circ\mathrm{C}$ là bao nhiêu $\mathrm{dm^3}$?
\shortans{44,8}
\loigiai{
Ta có
$1 \mathrm{atm}=760 \mathrm{mmHg}.$
Vì áp suất không đổi nên quá trình là đẳng áp.

Đổi nhiệt độ sang Kelvin:
$T_1=0+273=273 \mathrm{K},\qquad T_2=273+273=546 \mathrm{K}.$ Áp dụng định luật Charles:
 $\dfrac{V_1}{T_1}=\dfrac{V_2}{T_2}\Rightarrow$ V_2= $\dfrac{V_1T_2}{T_1}$. 
Thay số:V_2= $\dfrac{22,4\cdot 546}{273}=44,8\mathrm{dm^3}$.
Vậy thể tích của lượng khí là $44,8$ \mathrm{dm^3}.
}
\end{ex}

% ===== Câu 194 =====
% ID: L12C2B10-97-TLN
% Mức: VD
\begin{ex}
% ID: L12C2B10-97-TLN
% Mức: VD
Có $1\,\mathrm{g}$ khí hydrogen được đựng trong bình có thể tích $4\,\mathrm{L}$. Mật độ phân tử của chất khí đó là
\shortans{B}
\loigiai{
Ta có $n=\dfrac{m}{M}=\dfrac{1}{2}=0,5\ \mathrm{mol}.$ Số phân tử khí là $N=nN_A=0,5\cdot 6,02\cdot 10^{23}=3,01\cdot 10^{23}.$ Mật độ phân tử: $\mu=\dfrac{N}{V}=\dfrac{3,01\cdot 10^{23}}{4\cdot 10^{-3}}\approx 7,5\cdot 10^{25}\,\mathrm{m}^{-3}.$
}
\end{ex}

% ===== Câu 196 =====
% ID: L12C2B10-98-DS
% Mức: VD
\begin{ex}
% ID: L12C2B10-98-DS
% Mức: VD
Một bình dung tích $7,5\,\mathrm{L}$ chứa $24\,\mathrm{g}$ khí oxygen ở áp suất $2,5\cdot 10^5\,\mathrm{N/m^2}$.
\choiceTF
{Số mol khí oxygen trong bình là $0,5\,\mathrm{mol}$.}
{Số phân tử khí oxygen trong bình là $45\cdot 10^{23}$ phân tử.}
{\True Mật độ phân tử khí oxygen trong bình là $6,02\cdot 10^{25}\,\mathrm{m^{-3}}$.}
{\True Động năng phân tử trung bình của khí oxygen trong bình là $6,1\cdot 10^{-21}\,\mathrm{J}$.}
\loigiai{
\begin{itemchoice}
\item (Sai) Số mol khí oxygen trong bình là $0,5\,\mathrm{mol}$. (Vì): a sai. $n=\dfrac{m}{M}=\dfrac{24}{32}=0,75\,\mathrm{mol}.$ B. b sai. $N=nN_A=0,75\cdot 6,02\cdot 10^{23}=4,515\cdot 10^{23}\ \text{phân tử}.$ C. c đúng. $\mu=\dfrac{N}{V}=\dfrac{4,515\cdot 10^{23}}{7,5\cdot 10^{-3}}=6,02\cdot 10^{25}\,\mathrm{m^{-3}}.$ D. d đúng. Từ $p=\dfrac{2}{3}\mu\overline{W}_{\mathrm{đ}}$, suy ra $\overline{W}_{\mathrm{đ}}=\dfrac{3p}{2\mu}=\dfrac{3\cdot 2,5\cdot 10^5}{2\cdot 6,02\cdot 10^{25}}\approx6,2\cdot 10^{-21}\,\mathrm{J}.$
\item (Sai) Số phân tử khí oxygen trong bình là $45\cdot 10^{23}$ phân tử.
\item (Đúng) Mật độ phân tử khí oxygen trong bình là $6,02\cdot 10^{25}\,\mathrm{m^{-3}}$.
\item (Đúng) Động năng phân tử trung bình của khí oxygen trong bình là $6,1\cdot 10^{-21}\,\mathrm{J}$.
\end{itemchoice}
}
\end{ex}

% ===== Câu 198 =====
% ID: L12C2B10-99-TLN
% Mức: VD
\begin{ex}
% ID: L12C2B10-99-TLN
% Mức: VD
Tốc độ căn quân phương của khí có khối lượng riêng $2\,\mathrm{kg/m^3}$ ở áp suất $760\,\mathrm{mmHg}$ là bao nhiêu $\mathrm{m/s}$?
\shortans{390}
\loigiai{
Ta có $p=\dfrac{1}{3}\rho v^2\Rightarrow v=\sqrt{\dfrac{3p}{\rho}}.$ Với $p=760\,\mathrm{mmHg}=1,013\cdot 10^5\,\mathrm{Pa}$, ta được $v=\sqrt{\dfrac{3\cdot 1,013\cdot 10^5}{2}}\approx390\,\mathrm{m/s}.$
}
\end{ex}

% ===== Câu 200 =====
% ID: L12C2B10-100-TLN
% Mức: VD
\begin{ex}
% ID: L12C2B10-100-TLN
% Mức: VD
Nếu áp suất của một lượng khí biến đổi một lượng $2 \cdot 10^5 \text{ N/m}^2$ thì thể tích biến đổi một lượng là $3 \text{ lít}$. nếu áp suất biến đổi một lượng $5 \cdot 10^5 \text{ N/m}^2$ thì thể tích biến đổi một lượng là $3 \text{ lít}$ Coi nhiệt độ là không đổi thì áp suất và thể tích ban đầu của khí là
\shortans{A}
\loigiai{
Theo định luật Boyle-Mariotte cho quá trình đẳng nhiệt, tích của áp suất và thể tích của một lượng khí không đổi là một hằng số: $pV = \text{const}$.
 Gọi áp suất ban đầu của khí là $p$ (Pa) và thể tích ban đầu của khí là $V$ (lít).
 
 Khi áp suất biến đổi một lượng $2 \cdot 10^5 \text{ N/m}^2$ thì thể tích biến đổi một lượng $3 \text{ lít}$. Vì nhiệt độ không đổi, nếu áp suất tăng thì thể tích phải giảm để tích $pV$ không đổi.
 Trường hợp 1: Áp suất mới là $p_1' = p + 2 \cdot 10^5 \text{ Pa}$, thể tích mới là $V_1' = V - 3 \text{ lít}$.
 Áp dụng định luật Boyle:
 $pV = (p + 2 \cdot 10^5)(V - 3)pV = pV - 3p + 2 \cdot 10^5 V - (2 \cdot 10^5) \cdot 30 = -3p + 2 \cdot 10^5 V - 6 \cdot 10^53p - 2 \cdot 10^5 V = -6 \cdot 10^5$ (Phương trình 1)
 
 Khi áp suất biến đổi một lượng $5 \cdot 10^5 \text{ N/m}^2$ thì thể tích biến đổi một lượng $5 \text{ lít}$.
 Trường hợp 2: Áp suất mới là $p_2' = p + 5 \cdot 10^5 \text{ Pa}$, thể tích mới là $V_2' = V - 5 \text{ lít}$.
 Áp dụng định luật Boyle:
 $pV = (p + 5 \cdot 10^5)(V - 5)pV = pV - 5p + 5 \cdot 10^5 V - (5 \cdot 10^5) \cdot 50 = -5p + 5 \cdot 10^5 V - 25 \cdot 10^55p - 5 \cdot 10^5 V = -25 \cdot 10^5$
 Chia cả hai vế cho 5 để đơn giản hóa:
 $p - 1 \cdot 10^5 V = -5 \cdot 10^5$ (Phương trình 2)
 
 Ta có hệ phương trình gồm Phương trình 1 và Phương trình 2:
 $\begin{cases} 3p - 2 \cdot 10^5 V = -6 \cdot 10^5 \\ p - 1 \cdot 10^5 V = -5 \cdot 10^5 \end{cases}$
 
 Từ Phương trình 2, ta rút $p$ theo $V$:
 $p = 1 \cdot 10^5 V - 5 \cdot 10^5$ (Phương trình 3)
 
 Thay Phương trình 3 vào Phương trình 1:
 $3(1 \cdot 10^5 V - 5 \cdot 10^5) - 2 \cdot 10^5 V = -6 \cdot 10^53 \cdot 10^5 V - 15 \cdot 10^5 - 2 \cdot 10^5 V = -6 \cdot 10^5(3 - 2) \cdot 10^5 V = -6 \cdot 10^5 + 15 \cdot 10^51 \cdot 10^5 V = 9 \cdot 10^5V = \frac{9 \cdot 10^5}{1 \cdot 10^5}V = 9 \text{ lít}$
 
 Thay giá trị $V = 9 \text{ lít}$ vào Phương trình 3 để tìm $p$:
 $p = 1 \cdot 10^5 (9) - 5 \cdot 10^5p = 9 \cdot 10^5 - 5 \cdot 10^5p = 4 \cdot 10^5 \text{ Pa}$
 
 Vậy, áp suất ban đầu của khí là $p = 4 \cdot 10^5 \text{ Pa}$ và thể tích ban đầu của khí là $V = 9 \text{ lít}$.
 
 **A.** $V=9l\acute{i}t;p={{4.10}^{5}Pa.$
 Phương án này có giá trị thể tích $V = 9 \text{ lít}$ và áp suất $p = 4 \cdot 10^5 \text{ Pa}$, khớp hoàn toàn với kết quả tính toán của chúng ta. Do đó, đây là đáp án đúng.
 
 **B.** $V=9l\acute{i}t;p={{4.10}^{7}Pa.$
 Phương án này có thể tích $V = 9 \text{ lít}$ đúng với kết quả tính toán, nhưng áp suất $p = 4 \cdot 10^7 \text{ Pa}$ là sai, vì kết quả tính toán cho thấy áp suất ban đầu là $p = 4 \cdot 10^5 \text{ Pa}$.
 
 **C.** $V=9,5l\acute{i}t;p={{4.10}^{7}Pa.$
 Phương án này có thể tích $V = 9,5 \text{ lít}$ sai so với kết quả tính toán ($V = 9 \text{ lít}$), và áp suất $p = 4 \cdot 10^7 \text{ Pa}$ cũng sai so với kết quả tính toán ($p = 4 \cdot 10^5 \text{ Pa}$).
 
 **D.** $V=9,5l\acute{i}t;p={{4.10}^{5}Pa.$
 Phương án này có áp suất $p = 4 \cdot 10^5 \text{ Pa}$ đúng với kết quả tính toán, nhưng thể tích $V = 9,5 \text{ lít}$ là sai, vì kết quả tính toán cho thấy thể tích ban đầu là $V = 9 \text{ lít}$.
}
\end{ex}

% ===== Câu 202 =====
% ID: L12C2B10-101-TLN
% Mức: VD
\begin{ex}
% ID: L12C2B10-101-TLN
% Mức: VD
Một lượng không khí có thể tích 240 $c{{m}^{3}}$ chứa trong một xilanh có pit-tông đóng kín, diện tích của pit-tông là 24 $c{{m}^{2}}$ (xem hình vẽ bên). Áp suất của không khí trong xilanh bằng áp suất ngoài là $100\,\text{kPa}$. Cần một lực bằng bao nhiêu để dịch chuyển pit-tông $2\,\text{cm}$ theo chiều làm thể tích khí giảm? Bỏ qua ma sát giữa pit-tông và thành xilanh. Coi trong quá trình chuyển động nhiệt độ không thay đổi.
\shortans{B}
\loigiai{
* Lượng thể tích khí giảm: $\Delta V=Sh=24.2=48\left( c{{m}^{3}} \right)$ * Từ: ${{p}_{1}}{{V}_{1}}={{p}_{2}}{{V}_{2}}\xrightarrow[{}]{{{V}_{2}}={{V}_{1}}-\Delta V}{{p}_{2}}={{p}_{1}}\dfrac{{{V}_{1}}}{{{V}_{1}}-\Delta V}={{10}^{5}}.\dfrac{240}{240-48}=1,{{25.10}^{5}}\left( Pa \right)$ + Điều kiện pitông cân bằng: ${{p}_{t}}S={{p}_{n}}S+F\Rightarrow F=\left( {{p}_{t}}-{{p}_{n}} \right)S=\left( 1,{{25.10}^{5}}-{{10}^{5}} \right){{24.10}^{-4}}=60\left( N \right)$
}
\end{ex}

% ===== Câu 204 =====
% ID: L12C2B10-102-TLN
% Mức: VD
\begin{ex}
% ID: L12C2B10-102-TLN
% Mức: VD
Người ta bơm không khí áp suất 1 atm, vào bình có dung tích 10 lít. Biết mỗi lần bơm, bơm được 250 $c{{m}^{3}}$ không khí. Trước khi bơm đã có không khí 1 atm trong bình và trong khi bơm nhiệt độ không khí không đổi. Tính áp suất khí trong bình sau 50 lần bơm.
\shortans{A}
\loigiai{
:*Thể tích không khí ở áp suất 1 atm: ${{V}_{1}}=NV+{{V}_{2}}={{50.250.10}^{-3}}+10=22,5\left( l\acute{i}t \right)$ Theo định luật Boyle: ${{p}_{1}}{{V}_{1}}={{p}_{2}}{{V}_{2}}{{p}_{2}}={{p}_{1}}\dfrac{{{V}_{1}}}{{{V}_{2}}}=1.\dfrac{22,5}{10}=2,25\left( atm \right)$
}
\end{ex}

% ===== Câu 206 =====
% ID: L12C2B10-103-TLN
% Mức: VD
\begin{ex}
% ID: L12C2B10-103-TLN
% Mức: VD
Người ta dùng một bơm tay có ống bơm dài $50\,\text{cm}$ và đường kính trong $4\,\text{cm}$ để bơm khí vào một cái túi cao su sao cho túi phồng lên, có thể tích là 6,28 lít và áp suất của không khí trong túi là 4 atm. Biết áp suất của khí quyển là 1 atm và coi nhiệt độ của không khí được bơm vào túi không đổi. Số lần bơm đã thực hiện là bao nhiêu?
\shortans{40}
\loigiai{
Mỗi lần bơm, người ta đưa được vào túi cao su một lượng không khí có thể tích $V_{0}=\dfrac{\pi d^{2}\ell}{4}=\dfrac{3,14\cdot 4^{2}\cdot 50}{4}=628\,\mathrm{cm^{3}}$, áp suất $p_{0}=1\,\mathrm{atm}$. Khi được bơm vào túi ở áp suất $p=4\,\mathrm{atm}$, lượng khí này có thể tích $V$. Vì nhiệt độ không đổi nên theo định luật Boyle ta có: $pV=p_{0}V_{0}\Rightarrow V=\dfrac{p_{0}V_{0}}{p}=\dfrac{1\cdot 628}{4}=157\,\mathrm{cm^{3}}$. Số lần bơm: $n=\dfrac{6,28\cdot 10^{3}}{157}=40$.
}
\end{ex}

% ===== Câu 208 =====
% ID: L12C2B10-104-TLN
% Mức: VD
\begin{ex}
% ID: L12C2B10-104-TLN
% Mức: VD
Hai đường đẳng nhiệt của cùng một lượng khí lí tưởng ứng với hai nhiệt độ $T_1 = 300\ \mathrm{K}$ và $T_2 = 450\ \mathrm{K}$ được vẽ trên đồ thị $(p, V)$. Hãy tính tỉ số giữa tích $p \cdot V$ tại một điểm bất kì trên đường $T_2$ so với đường $T_1$.
\shortans{1,5}
\loigiai{
Áp dụng phương trình liên hệ thông số trạng thái vĩ mô của lượng khí lí tưởng xác định:
  \begin{aligned}
 $\dfrac{p_2 \cdot V_2}{p_1 \cdot V_1}&=\dfrac{T_2}{T_1}&=\dfrac{450}{300}$ 
&= 1,5.
 \end{aligned} 
 Kết luận: Tỉ số giữa hai tích bằng $1,5$.
}
\end{ex}

% ===== Câu 209 =====
% ID: L12C2B10-105-TN
% Mức: VD
\begin{ex}
% ID: L12C2B10-105-TN
% Mức: VD
Một lượng khí đựng trong xilanh có pit-tông chuyển động được. Các thông số trạng thái của lượng khí này là: $2\ \text{atm}$, $15\ \text{lít}$, $300\ \text{K}$. Khi pit-tông nén khí, áp suất của khí tăng lên tới $3,5\ \text{atm}$, thể tích giảm còn $12\ \text{lít}$. Nhiệt độ của khí nén bằng
\choice
{\True $420\ \text{K}$}
{$360\ \text{K}$}
{$480\ \text{K}$}
{$257\ \text{K}$}
\loigiai{
Áp dụng phương trình trạng thái khí lí tưởng cho hai trạng thái của khí:
$\frac{p_1V_1}{T_1}=\frac{p_2V_2}{T_2}$

Thay số liệu:
- Trạng thái 1: $p_1 = 2\ \text{atm}$, $V_1 = 15\ \text{lít}$, $T_1 = 300\ \text{K}$
- Trạng thái 2: $p_2 = 3,5\ \text{atm}$, $V_2 = 12\ \text{lít}$, $T_2 = ?$ \frac{2 \times 15}{300}=\frac{3,5 \times 12}{T_2} $\frac{30}{300}=\frac{42}{T_2}$ T_2 = \frac{42 \times 300}{30} = \frac{12600}{30} = 420\ \text{K} $Vậy nhiệt độ của khí nén là$ 420\ \text{K}.
}
\end{ex}

% ===== Câu 210 =====
% ID: L12C2B10-106-TLN
% Mức: VD
\begin{ex}
% ID: L12C2B10-106-TLN
% Mức: VD
Nếu thể tích một lượng khí giảm $\dfrac{1}{10}$ thì áp suất tăng $\dfrac{1}{5}$ so với áp suất ban đầu và nhiệt độ tăng thêm $16^\circ\mathrm{C}$. Nhiệt độ ban đầu của khối khí là
\shortans{A}
\loigiai{
Theo đề bài:
$V_2=\dfrac{9}{10}V_1,\qquad p_2=\dfrac{6}{5}p_1,\qquad T_2=T_1+16.$

Áp dụng phương trình trạng thái khí lí tưởng:
$\dfrac{p_1V_1}{T_1}=\dfrac{p_2V_2}{T_2}.$

Suy ra
$\dfrac{p_1V_1}{T_1} = \dfrac{\dfrac{6}{5}p_1\cdot \dfrac{9}{10}V_1}{T_1+16}.$

Rút gọn: 
$\dfrac{1}{T_1}=\dfrac{1,08}{T_1+16}.$

Do đó
$T_1+16=1,08T_1\Rightarrow$ 0,08T_1=16 $\Rightarrow T_1=200 \text{ K}.$
}
\end{ex}

% ===== Câu 211 =====
% ID: L12C2B10-107-TN
% Mức: VD
\begin{ex}
% ID: L12C2B10-107-TN
% Mức: VD
Hình dưới là đường đẳng áp của cùng một lượng khí ứng với các áp suất $p_1$, $p_2$ và $p_3$. Sắp xếp nào dưới đây là đúng?
\choice
{\True $p_1 \lt  p_2 \lt  p_3$}
{$p_1 \gt  p_2 \gt  p_3$}
{$p_2 \gt  p_1 \gt  p_3$}
{$p_2 \lt  p_1 \lt  p_3$}
\loigiai{
Áp suất khí trong mô hình động học phân tử liên quan đến thể tích và nhiệt độ theo phương trình trạng thái khí lí tưởng: $pV = nRT$. Đường đẳng áp là đường có áp suất không đổi, nhưng đề bài cho hình có các áp suất khác nhau $p_1, p_2, p_3$ trên cùng một đường đẳng áp, tức là cần so sánh các áp suất dựa trên vị trí trên đồ thị. 
 A. $p_1 \lt  p_2 \lt  p_3$: Nếu thể tích giảm dần từ $V_1 \gt  V_2 \gt  V_3$ trên đường đẳng nhiệt, áp suất tăng dần, nên đúng. 
 B. $p_1 \gt  p_2 \gt  p_3$: Ngược lại với quan hệ thể tích và áp suất trên đường đẳng nhiệt, sai. 
 C. $p_2 \gt  p_1 \gt  p_3$: Không đúng vì thứ tự áp suất không tương ứng với thứ tự thể tích. 
 D. $p_2 \lt  p_1 \lt  p_3$: Không đúng vì không phù hợp với mối quan hệ thể tích và áp suất trên đường đẳng nhiệt.
}
\end{ex}

% ===== Câu 212 =====
% ID: L12C2B10-108-TLN
% Mức: VD
\begin{ex}
% ID: L12C2B10-108-TLN
% Mức: VD
Một mol khí ở áp suất $2\,\mathrm{atm}$ và nhiệt độ $30^\circ\mathrm{C}$ thì chiếm một thể tích là
\shortans{C}
\loigiai{
Ở điều kiện tiêu chuẩn, một mol khí chiếm thể tích $22,4$ lít.

Ta có
$T=30+273=303 \mathrm{K}.$
Áp dụng phương trình trạng thái:
$\dfrac{pV}{T}=\text{hằng số}.$
Suy ra
 V= $\dfrac{1\cdot 22,4\cdot 303}{2\cdot 273}\approx$ 12,43\text{ lít}.
}
\end{ex}

% ===== Câu 214 =====
% ID: L12C2B10-109-TLN
% Mức: VD
\begin{ex}
% ID: L12C2B10-109-TLN
% Mức: VD
Một lượng khí có áp suất $750\,\mathrm{mmHg}$, nhiệt độ $27^\circ\mathrm{C}$ và thể tích $76\,\mathrm{cm^3}$. Thể tích khí ở điều kiện chuẩn $\left(0^\circ\mathrm{C},\,760\,\mathrm{mmHg}\right)$ là
\shortans{C}
\loigiai{
Ta có
$T_1=27+273=300 \mathrm{K},\qquad T_0=273 \mathrm{K}.$
Áp dụng phương trình trạng thái khí lí tưởng:
$\dfrac{p_1V_1}{T_1}=\dfrac{p_0V_0}{T_0}.$
Suy ra
 V_0= $\dfrac{p_1V_1T_0}{T_1p_0}$ = $\dfrac{750\cdot 76\cdot 273}{300\cdot 760}\approx$ 68,25 $\mathrm{cm^3}$.
}
\end{ex}

% ===== Câu 216 =====
% ID: L12C2B10-110-TLN
% Mức: VD
\begin{ex}
% ID: L12C2B10-110-TLN
% Mức: VD
Ở $27^\circ\mathrm{C}$, thể tích của một lượng khí là $6$ lít. Thể tích của lượng khí đó ở nhiệt độ $227^\circ\mathrm{C}$ khi áp suất không đổi là
\shortans{B}
\loigiai{
Ta có
$T_1=27+273=300 \mathrm{K},\qquad V_1=6\text{ lít}.$
Ở trạng thái sau:
$T_2=227+273=500 \mathrm{K}.$
Quá trình biến đổi là đẳng áp nên
$\dfrac{V_1}{T_1}=\dfrac{V_2}{T_2}.$
Suy ra
 V_2= $\dfrac{V_1T_2}{T_1}$ =\dfrac{6$\cdot 500}{300}=10\text{ lít}.
}
\end{ex}

% ===== Câu 218 =====
% ID: L12C2B10-111-TLN
% Mức: VD
\begin{ex}
% ID: L12C2B10-111-TLN
% Mức: VD
Một khối khí có thể tích $1\,\mathrm{m^3}$, nhiệt độ $11^\circ\mathrm{C}$. Để giảm thể tích khí còn một nửa khi áp suất không đổi cần
\shortans{C}
\loigiai{
Ta có
$T_1=11+273=284 \mathrm{K}.$
Do áp suất không đổi nên
$\dfrac{V_1}{T_1}=\dfrac{V_2}{T_2}.$
Khi thể tích giảm còn một nửa:
$V_2=\dfrac{V_1}{2}.$ Suy ra
 T_2= $\dfrac{V_2}{V_1}T_1$ = $\dfrac{1}{2}\cdot$ 284=142 $\mathrm{K}$. 
Đổi sang độ Celsius:t_2=142-273=-131^ $\circ\mathrm{C}$.
}
\end{ex}

% ===== Câu 220 =====
% ID: L12C2B10-112-TLN
% Mức: VD
\begin{ex}
% ID: L12C2B10-112-TLN
% Mức: VD
Trong xi lanh của một động cơ đốt trong có $40$ lít hỗn hợp khí đốt dưới áp suất $1\,\mathrm{atm}$ và nhiệt độ $47^\circ\mathrm{C}$. Pít-tông nén xuống làm cho thể tích của hỗn hợp khí chỉ còn $5$ lít và áp suất tăng lên đến $15\,\mathrm{atm}$. Nhiệt độ hỗn hợp khí khi đó là
\shortans{B}
\loigiai{
Ta có
$T_1=47+273=320 \mathrm{K}.$
Áp dụng phương trình trạng thái khí lí tưởng:
$\dfrac{p_1V_1}{T_1}=\dfrac{p_2V_2}{T_2}.$ Suy ra
 T_2= $\dfrac{p_2V_2T_1}{p_1V_1}$ = $\dfrac{15\cdot 5\cdot 320}{1\cdot 40}$ =600 $\mathrm{K}$. 
Vậyt_2=600-273=327^ $\circ\mathrm{C}$.
}
\end{ex}

% ===== Câu 222 =====
% ID: L12C2B10-113-TLN
% Mức: VD
\begin{ex}
% ID: L12C2B10-113-TLN
% Mức: VD
Khối lượng riêng của không khí ở nhiệt độ $80^\circ\mathrm{C}$ và có áp suất $2,5\cdot 10^5\,\mathrm{Pa}$. Biết khối lượng riêng của không khí ở $0^\circ\mathrm{C}$ là $1,29\,\mathrm{kg/m^3}$ và áp suất $1,01\cdot 10^5\,\mathrm{Pa}$ là
\shortans{C}
\loigiai{
Với cùng một lượng khí:
$\dfrac{pV}{T}=\text{hằng số}.$
Mà
$m=\rho V \Rightarrow V=\dfrac{m}{\rho}.$
Suy ra
$\dfrac{p}{\rho T}=\text{hằng số}.$
Do đó
$\rho_2=\dfrac{\rho_1T_1p_2}{T_2p_1}.$
Thay số:
 $\rho_2$ = $\dfrac{1,29\cdot 273\cdot 2,5\cdot 10^5}$ {353 $\cdot$ 1,01 $\cdot$ 10^5}
$\approx$ 2,47 $\mathrm{kg/m^3}$.
}
\end{ex}

% ===== Câu 224 =====
% ID: L12C2B10-114-TLN
% Mức: VD
\begin{ex}
% ID: L12C2B10-114-TLN
% Mức: VD
Nén $10$ lít khí ở nhiệt độ $27^\circ\mathrm{C}$ để cho thể tích của nó chỉ còn $4$ lít, vì nén nhanh nên khí bị nóng lên đến $60^\circ\mathrm{C}$. Áp suất chất khí tăng lên
\shortans{B}
\loigiai{
Ta có
$T_1=27+273=300 \mathrm{K},\qquad T_2=60+273=333 \mathrm{K}.$
Áp dụng phương trình trạng thái khí lí tưởng:
$\dfrac{p_1V_1}{T_1}=\dfrac{p_2V_2}{T_2}.$
Suy ra
 $\dfrac{p_2}{p_1}$ = $\dfrac{V_1}{V_2}\cdot\dfrac{T_2}{T_1}$ = $\dfrac{10}{4}\cdot\dfrac{333}{300}\approx$ 2,78. 
Vậy áp suất chất khí tăng lên khoảng $2,78$ lần.
}
\end{ex}

% ===== Câu 226 =====
% ID: L12C2B10-115-TLN
% Mức: VD
\begin{ex}
% ID: L12C2B10-115-TLN
% Mức: VD
Khí trong một bình dung tích $3\,\mathrm{dm^3}$, áp suất $200\,\mathrm{kPa}$ và nhiệt độ $16^\circ\mathrm{C}$ có khối lượng $11$ gam. Khí đó là
\shortans{C}
\loigiai{
Đổi:
$p=200 \mathrm{kPa}=2\cdot 10^5 \mathrm{Pa},V=3 \mathrm{dm^3}=3\cdot 10^{-3} \mathrm{m^3},T=16+273=289 \mathrm{K},\qquad m=11 \mathrm{g}.$
Từ phương trình khí lí tưởng:
$pV=nRT.$
Suy ra
$n=\dfrac{pV}{RT}.$
Khối lượng mol:
$M=\dfrac{m}{n}=\dfrac{mRT}{pV}.$
Thay số:
 M= $\dfrac{11\cdot 8,31\cdot 289}{2\cdot 10^5\cdot 3\cdot 10^{-3}}\approx$ 44 $\mathrm{g/mol}$. 
Vậy khí đó là $\mathrm{CO_2}$.
}
\end{ex}

% ===== Câu 228 =====
% ID: L12C2B10-116-TLN
% Mức: VD
\begin{ex}
% ID: L12C2B10-116-TLN
% Mức: VD
Trong xi lanh của một động cơ đốt trong có $2\,\mathrm{dm^3}$ hỗn hợp khí đốt dưới áp suất $1\,\mathrm{atm}$ và nhiệt độ $27^\circ\mathrm{C}$. Pít-tông nén xuống làm cho thể tích của hỗn hợp khí chỉ còn $0,2\,\mathrm{dm^3}$ và áp suất tăng thêm $15\,\mathrm{atm}$. Nhiệt độ hỗn hợp khí khi đó là
\shortans{A}
\loigiai{
Ta có
$T_1=27+273=300 \mathrm{K}.$
Vì áp suất tăng thêm $15\,\mathrm{atm}$ nên
$p_2=1+15=16 \mathrm{atm}.$
Áp dụng phương trình trạng thái khí lí tưởng:
$\dfrac{p_1V_1}{T_1}=\dfrac{p_2V_2}{T_2}.$ Suy ra
 T_2= $\dfrac{p_2V_2T_1}{p_1V_1}$ = $\dfrac{16\cdot 0,2\cdot 300}{1\cdot 2}$ =480 $\mathrm{K}$. 
Vậyt_2=480-273=207^ $\circ\mathrm{C}$.
}
\end{ex}

% ===== Câu 230 =====
% ID: L12C2B10-117-TN
% Mức: VDC
\begin{ex}
% ID: L12C2B10-117-TN
% Mức: VDC
Pít-tông của một máy nén sau mỗi lần nén đưa được $4$ lít khí ở nhiệt độ $27^\circ\mathrm{C}$ và áp suất $1\,\mathrm{atm}$ vào bình chứa khí có thể tích $3\,\mathrm{m^3}$. Khi pít-tông đã thực hiện $1000$ lần nén và nhiệt độ khí trong bình là $42^\circ\mathrm{C}$ thì áp suất khí trong bình nhận giá trị là
\choice
{$1,79\,\mathrm{atm}$}
{$1,27\,\mathrm{atm}$}
{\True $2,45\,\mathrm{atm}$}
{$2,9\,\mathrm{atm}$}
\loigiai{
Đổi:
$3 \mathrm{m^3}=3000\text{ lít}.$
Sau $1000$ lần nén, tổng thể tích khí quy về điều kiện ban đầu là
$V_1=4\cdot 1000+3000=7000\text{ lít}.$
Ta có
$T_1=27+273=300 \mathrm{K},\qquad T_2=42+273=315 \mathrm{K}.$
Áp dụng phương trình trạng thái khí lí tưởng:
$\dfrac{p_1V_1}{T_1}=\dfrac{p_2V_2}{T_2}.$
Suy ra
 p_2= $\dfrac{p_1V_1T_2}{T_1V_2}$ = $\dfrac{1\cdot 7000\cdot 315}{300\cdot 3000}$ =2,45 $\mathrm{atm}$.
}
\end{ex}

% ===== Câu 232 =====
% ID: L12C2B10-118-TN
% Mức: VDC
\begin{ex}
% ID: L12C2B10-118-TN
% Mức: VDC
Coi áp suất khí trong và ngoài phòng như nhau. Khối lượng riêng của không khí trong phòng ở nhiệt độ $27^\circ\mathrm{C}$ lớn hơn khối lượng riêng của không khí ngoài sân nắng ở nhiệt độ $42^\circ\mathrm{C}$ bao nhiêu lần?
\choice
{$1{,}5$ lần}
{\True $1{,}05$ lần}
{$10{,}5$ lần}
{$15$ lần}
\loigiai{
Với cùng một chất khí và cùng áp suất, khối lượng riêng tỉ lệ nghịch với nhiệt độ tuyệt đối:
$D\sim \dfrac{1}{T}.$
Ta có
$T_1=27+273=300 \mathrm{K},\qquad T_2=42+273=315 \mathrm{K}.$
Do đó
 $\dfrac{D_1}{D_2}=\dfrac{T_2}{T_1}$ = $\dfrac{315}{300}=1,05.$ 
Vậy khối lượng riêng của không khí trong phòng lớn hơn khối lượng riêng của không khí ngoài sân nắng $1,05$ lần.
}
\end{ex}

% ===== Câu 234 =====
% ID: L12C2B10-119-TN
% Mức: VDC
\begin{ex}
% ID: L12C2B10-119-TN
% Mức: VDC
Hai bình cầu được nối với nhau bằng một ống có khóa, chứa hai chất khí không tác dụng hóa học với nhau, ở cùng nhiệt độ. Áp suất khí trong hai bình lần lượt là $p_1=2\cdot 10^5\,\mathrm{N/m^2}$ và $p_2=10^6\,\mathrm{N/m^2}$. Mở khóa nhẹ nhàng để hai bình thông với nhau sao cho nhiệt độ không đổi. Khi cân bằng xảy ra, áp suất ở hai bình là $p=4\cdot 10^5\,\mathrm{N/m^2}$. Tỉ số thể tích của hai bình cầu là
\choice
{\True $3$}
{$4$}
{$5$}
{$1$}
\loigiai{
Gọi
$n=\dfrac{V_1}{V_2}.$
Khi mở khóa, nhiệt độ không đổi nên áp suất cuối cùng của hỗn hợp khí là
$p=\dfrac{p_1V_1+p_2V_2}{V_1+V_2}.$
Chia cả tử và mẫu cho $V_2$, ta được
$p=\dfrac{p_1n+p_2}{n+1}= 4 \cdot 10^5$ = $\dfrac{2\cdot 10^5\cdot n+10^6}{n+1}$. 
Chia hai vế cho $10^5$:
$4=\dfrac{2n+10}{n+1}.$ Suy ra
 $4\,\text{n}$ +4= $2\,\text{n}$ +10
$\Rightarrow2\,\text{n}$ =6
$\Rightarrow$ n=3. 
Vậy $\dfrac{V_1}{V_2}=3.$
}
\end{ex}

% ===== Câu 236 =====
% ID: L12C2B10-120-TN
% Mức: VDC
\begin{ex}
% ID: L12C2B10-120-TN
% Mức: VDC
Có $12$ gam khí chiếm thể tích $4$ lít ở $7^\circ\mathrm{C}$. Sau khi nung nóng đẳng áp, khối lượng riêng của khí là $1,2\,\mathrm{g/l}$. Nhiệt độ khí sau khi nung là
\choice
{\True $427^\circ\mathrm{C}$}
{$372^\circ\mathrm{C}$}
{$327\,\mathrm{K}$}
{$372\,\mathrm{K}$}
\loigiai{
Thể tích khí sau khi nung là
$V_2=\dfrac{m}{\rho_2}=\dfrac{12}{1,2}=10\text{ lít}.$
Ta có
$T_1=7+273=280 \mathrm{K}.$ Áp dụng định luật Charles:
 $\dfrac{V_1}{T_1}=\dfrac{V_2}{T_2}\Rightarrow$ T_2=T_1 $\dfrac{V_2}{V_1}$ =280 $\cdot\dfrac{10}{4}$ =700 $\mathrm{K}$. 
Suy rat_2=T_2-273=700-273=427^ $\circ\mathrm{C}$.
}
\end{ex}

% ===== Câu 238 =====
% ID: L12C2B10-121-DS
% Mức: VDC
\begin{ex}
% ID: L12C2B10-121-DS
% Mức: VDC
Một căn phòng có thể tích $V=60\,\mathrm{m^3}$. Khi ta tăng nhiệt độ của phòng từ $T_1=280\,\mathrm{K}$ đến $T_2=300\,\mathrm{K}$ ở áp suất $101,3\,\mathrm{kPa}$. Cho biết khối lượng riêng của không khí ở điều kiện chuẩn, nhiệt độ $T_0=273\,\mathrm{K}$, áp suất $p_0=101,3\,\mathrm{kPa}$, là $\rho_0=1,29\,\mathrm{kg/m^3}$.
\choiceTF
{\True Có thể áp dụng định luật Charles cho quá trình biến đổi trạng thái này}
{Khi tăng nhiệt độ đẳng áp thì độ tăng thể tích được xác định bởi biểu thức $\Delta V=\dfrac{T_1-T_2}{T_1}V$}
{\True Khối lượng riêng của không khí ở điều kiện sau khi tăng nhiệt độ xấp xỉ bằng $1,17\,\mathrm{kg/m^3}$}
{Khối lượng không khí thoát ra khỏi phòng là $5,013\,\mathrm{kg}$}
\loigiai{
\begin{itemchoice}
\item (Đúng) Có thể áp dụng định luật Charles cho quá trình biến đổi trạng thái này. (Vì): ] Đúng. Do áp suất không đổi nên có thể áp dụng định luật Charles cho quá trình biến đổi trạng thái này
\item (Sai) Khi tăng nhiệt độ đẳng áp thì độ tăng thể tích được xác định bởi biểu thức $\Delta V=\dfrac{T_1-T_2}{T_1}V$. (Vì): ] Sai.

 Xét lượng không khí ban đầu trong phòng. Khi tăng nhiệt độ đẳng áp, nếu không bị giới hạn bởi thể tích phòng thì lượng khí đó có thể tích mới là $V_2$.

 Theo định luật Charles:
 $\dfrac{V_2}{V}=\dfrac{T_2}{T_1}.$ Suy ra
  $\dfrac{\Delta V}{V}$ = $\dfrac{V_2-V}{V}$ = $\dfrac{T_2-T_1}{T_1}$. 
 Do đó $\Delta V=\dfrac{T_2-T_1}{T_1}V.$
 Vậy biểu thức $\Delta V=\dfrac{T_1-T_2}{T_1}V$ là sai
\item (Đúng) Khối lượng riêng của không khí ở điều kiện sau khi tăng nhiệt độ xấp xỉ bằng $1,17\,\mathrm{kg/m^3}$.
\item (Sai) Khối lượng không khí thoát ra khỏi phòng là $5,013\,\mathrm{kg}$.
\end{itemchoice}
}
\end{ex}

% ===== Câu 240 =====
% ID: L12C2B10-122-TLN
% Mức: VDC
\begin{ex}
% ID: L12C2B10-122-TLN
% Mức: VDC
Một áp kế gồm một bình cầu thủy tinh có thể tích $270\,\mathrm{cm^3}$ gắn với ống nhỏ $AB$ nằm ngang có tiết diện $0,1\,\mathrm{cm^2}$. Trong ống có một giọt thủy ngân. Ở $0^\circ\mathrm{C}$, giọt thủy ngân cách $A$ một đoạn $30\,\mathrm{cm}$. Hỏi khi nung bình đến $10^\circ\mathrm{C}$ thì giọt thủy ngân ở vị trí cách $A$ một khoảng bao nhiêu? Coi dung tích của bình không đổi, ống $AB$ đủ dài để giọt thủy ngân không chảy ra ngoài.
\shortans{130}
\loigiai{
Gọi $x$ là khoảng cách từ $A$ đến vị trí mới của giọt thủy ngân, tính bằng $\mathrm{cm}$.

Thể tích khí ban đầu là
$V_1=270+0,1\cdot 30=273 \mathrm{cm^3}.$
Thể tích khí sau khi nung là
$V_2=270+0,1x.$
Đổi nhiệt độ sang Kelvin:
$T_1=0+273=273 \mathrm{K},\qquad T_2=10+273=283 \mathrm{K}.$
Quá trình nung nóng là quá trình đẳng áp, nên:
$\dfrac{V_1}{T_1}=\dfrac{V_2}{T_2}.$
Suy ra
$\dfrac{273}{273}=\dfrac{270+0,1x}{283}.$
Do đó
 270+0,1x=283
$\Rightarrow$ 0,1x=13
$\Rightarrow$ x=130 $\mathrm{cm}$. 
Vậy giọt thủy ngân ở vị trí cách $A$ một khoảng $130\,\mathrm{cm}$.
}
\end{ex}

% ===== Câu 242 =====
% ID: L12C2B10-123-TLN
% Mức: VDC
\begin{ex}
% ID: L12C2B10-123-TLN
% Mức: VDC
Một xi lanh cách nhiệt hình trụ chiều cao h= $50\,\text{cm}$, tiết diện $S=100cm^{2}$ đặt thẳng đứng, xi lanh được chia thành hai phần bằng một pittong cách nhiệt khối lượng m= $500\,\text{g}$. Khí trong hai phần cùng loại ở nhiệt độ $({20)$ ^{0}}C $và có khối lượng$ ({m) $_{1}}$ =0,5 $({m)$ _{2}} $. Pittong cân bằng khi ở cách đáy dưới đoạn$ ({h) $_{2}}=0,4h.$ Tính áp suất khí trong phần dưới của xi lanh theo đơn vị kPa? Lấy $g=10m/s^{2}$
\shortans{2}
\loigiai{
Khi pitton đứng cân bằng thì: ${{F}_{2}}={{F}_{1}}+{{F}_{0}}\Leftrightarrow {{p}_{2}}={{p}_{1}}+\dfrac{mg}{S}$ (1) + Áp dụng PTTT cho từng lượng khí ta có ${{p}_{1}}{{V}_{1}}=\dfrac{m}{{{M}_{1}}}R{{T}_{1}}$ và ${{p}_{2}}.{{V}_{2}}=\dfrac{m}{{{M}_{2}}}R.{{T}_{2}}$ + Do ${{T}_{1}}={{T}_{2}}$ nên ta có $\dfrac{{{p}_{1}}{{V}_{1}}}{{{p}_{2}}.{{V}_{2}}}=\dfrac{1}{2}\Rightarrow \dfrac{{{p}_{1}}}{{{p}_{2}}}=\dfrac{3}{4}(2)$ Từ (1,2) ta có ${{p}_{2}}=2000Pa=2kPa$
}
\end{ex}

% ===== Câu 244 =====
% ID: L12C2B10-124-TLN
% Mức: VDC
\begin{ex}
% ID: L12C2B10-124-TLN
% Mức: VDC
Trong một ống nhỏ dài, một đầu kín, một đầu hở, có tiết diện đều. Ban đầu, đặt ống thẳng đứng sao cho miệng ống hướng lên trên. Trong ống, về phía đáy có một cột không khí dài $30\mathrm{cm}$, được ngăn cách với bên ngoài bởi một cột thủy ngân dài $h=15\mathrm{cm}$.

Biết áp suất khí quyển là $76\mathrm{cmHg}$ và nhiệt độ không đổi. Tính chiều dài cột không khí trong ống khi đặt ống nghiêng góc $30^\circ$ so với phương ngang, miệng ống ở phía trên.
\shortans{32,7}
\loigiai{
Cột thủy ngân có độ dài $h$. Khi đặt ống nghiêng góc $30^\circ$ so với phương ngang, độ cao của cột thủy ngân là
$h'=h\sin30^\circ=\dfrac{h}{2}=7,5 \mathrm{cm}.$


    
• Khi ống thẳng đứng, áp suất của khối khí là
    $p_1=p_0+h=76+15=91 \mathrm{cmHg}.$

    
• Khi ống đặt nghiêng, áp suất của khối khí là
    $p_3=p_0+h'=76+7,5=83,5 \mathrm{cmHg}.$

    
• Theo định luật Boyle:
    $p_1V_1=p_3V_3.$

    Với
    $V_1=30S,\qquad V_3=\ell_3S,$
    ta có
    $91\cdot 30S=83,5\cdot \ell_3S.$

    Suy ra
    $\ell_3=\dfrac{91\cdot 30}{83,5}\approx 32,7 \mathrm{cm}.$


Vậy chiều dài cột không khí khi ống đặt nghiêng là
$\boxed{\ell_3\approx 32,7 \mathrm{cm}}.$
}
\end{ex}

% ===== Câu 246 =====
% ID: L12C2B10-125-TLN
% Mức: VDC
\begin{ex}
% ID: L12C2B10-125-TLN
% Mức: VDC
Trong một ống hình trụ thẳng đứng có hai tiết diện khác nhau có hai pittong nối với nhau bằng sợi dây không dãn, giữa hai pittong có một mol khí lí tưởng. Pittong trên có tiết diện lớn hơn pittong dưới $\Delta S=10{cm}^{2}$. Áp suất khí quyển bên ngoài là 1 atm.Tính áp suất p của khí giữa hai pitton theo đơn vị atm). Biết khối lượng tổng cộng của hai pittong là m= $5\,\text{kg}$, $1 atm\approx {10}^{5}Pa$.
\shortans{1,5}
\loigiai{
Ta có khi pitong đứng cân bằng: 


• pittong 1: $p_{o.}S_{1}+m_{1}g=p.S_{1}$ (1) 


• pittong 2: $p_{.}S_{2}+m_{2}g=p_{0}.S_{2}$ (2)

 lấy (1) + (2) ta suy ra $P$ = $p_{0}+\dfrac{mg}{\Delta S}$ thay số p =1,5atm
}
\end{ex}

% ===== Câu 248 =====
% ID: L12C2B10-126-TLN
% Mức: VDC
\begin{ex}
% ID: L12C2B10-126-TLN
% Mức: VDC
Hai bình có thể tích lần lượt là $V_1=40$ lít, $V_2=10$ lít thông nhau qua một cái van. Van chỉ mở khi áp suất trong bình $1$ lớn hơn trong bình $2$ từ $10^5\,\mathrm{Pa}$ trở lên. Ban đầu bình $1$ chứa khí ở áp suất $p_0=0,9\cdot 10^5\,\mathrm{Pa}$ và nhiệt độ $T_0=300\,\mathrm{K}$, còn bình $2$ là chân không. Người ta làm nóng đều cả hai bình từ nhiệt độ $T_0$ lên nhiệt độ $T=500\,\mathrm{K}$. Áp suất cuối cùng trong bình $2$ là
\shortans{D}
\loigiai{
Van bắt đầu mở khi áp suất bình $1$ đạt
$p_m=10^5 \mathrm{Pa}.$ Nhiệt độ lúc van bắt đầu mở là
 T_m= $\dfrac{p_m}{p_0}T_0$ = $\dfrac{10^5}{0,9\cdot 10^5}\cdot$ 300
\approx 333 $\mathrm{K}$. 
Sau đó, van giữ cho hiệu áp suất giữa hai bình là $p_1-p_2=10^5 \mathrm{Pa}.$
Ở nhiệt độ cuối $T=500 \mathrm{K}$, bảo toàn lượng khí:
 \dfrac{p_0V_1}{T_0}
=
\dfrac{p_1V_1+p_2V_2}{T}. 
Với
$p_1=p_2+10^5,\qquad V_1=40,\qquad V_2=10.$ Thay vào:
 $\dfrac{0,9\cdot 10^5\cdot 40}{300}$ = $\dfrac{(p_2+10^5)\cdot 40+p_2\cdot 10}{500}$. 
Giải ra:p_2=0,4\cdot 10^5 $\mathrm{Pa}$.
}
\end{ex}

% ===== Câu 250 =====
% ID: L12C2B10-127-TLN
% Mức: VDC
\begin{ex}
% ID: L12C2B10-127-TLN
% Mức: VDC
Một ống thủy tinh hình trụ có tiết diện không đổi, một đầu kín, được dùng làm ống Torricelli để đo áp suất khí quyển. Vì có một ít không khí ở trong ống trên mức thủy ngân nên khi áp suất khí quyển là $p_0$, ở nhiệt độ $T_0$ thì chiều cao cột thủy ngân là $H_0$. Biết chiều dài của ống từ mặt thủy ngân trong chậu đến đầu trên được giữ không đổi và bằng $L$. Nếu ở nhiệt độ $T$ chiều cao cột thủy ngân là $H$ thì áp suất khí quyển $p_k$ là
\shortans{C}
\loigiai{
Gọi $p_1$ là áp suất của không khí trong ống ở nhiệt độ $T_0$, $p$ là áp suất của không khí trong ống ở nhiệt độ $T$.

Ta có
 p_0=p_1+H_0
$\Rightarrow$ p_1=p_0-H_0. 
Ở trạng thái sau:
$p_k=p+H.$
Áp dụng phương trình trạng thái cho lượng khí trong ống:
 $\dfrac{p_1(L-H_0)}{T_0}$ = $\dfrac{p(L-H)}{T}$. 
Suy ra
 p
=
$(p_0-H_0)\dfrac{L-H_0}{L-H}\cdot\dfrac{T}{T_0}$. 
Vậy
 p_k
=
H+ $(p_0-H_0)\dfrac{L-H_0}{L-H}\cdot\dfrac{T}{T_0}$.
}
\end{ex}

% ===== Câu 252 =====
% ID: L12C2B10-128-TN
% Mức: VDC
\begin{ex}
% ID: L12C2B10-128-TN
% Mức: VDC
Một xi lanh đặt nằm ngang. Lúc đầu pít-tông cách đều hai đầu xi lanh một khoảng $50\,\mathrm{cm}$ và không khí chứa trong xi lanh có nhiệt độ $27^\circ\mathrm{C}$, áp suất $1\,\mathrm{atm}$. Sau đó không khí ở đầu bên trái được nung lên đến $t^\circ\mathrm{C}$ thì pít-tông dịch chuyển một khoảng $x=3\,\mathrm{cm}$. Nhiệt độ nung $t^\circ\mathrm{C}$ bằng
\choice
{\True $65,3^\circ\mathrm{C}$}
{$56^\circ\mathrm{C}$}
{$75^\circ\mathrm{C}$}
{$57^\circ\mathrm{C}$}
\loigiai{
Ban đầu, pít-tông cách đều hai đầu xi lanh nên hai bên có cùng lượng khí.

Khi nung nóng khí bên trái, pít-tông dịch chuyển sang phải $3\,\mathrm{cm}$. Khi đó:
$l_{\text{trái}}=50+3=53 \mathrm{cm},$ l_{\text{phải}}=50-3=47 $\mathrm{cm}$. $Khi pít-tông cân bằng, áp suất hai bên bằng nhau:$ p_{\text{trái}}=p_{\text{phải}}. $Với cùng lượng khí, ta có$ p= $\dfrac{nRT}{V}$. $Do đó\dfrac{t+273}{53}=\dfrac{300}{47}$. $Suy ra$ t+273= $\dfrac{300\cdot 53}{47}\approx$ 338,3. $Vậy$ t\approx 65,3^ $\circ\mathrm{C}$.
}
\end{ex}
